% ═══════════════════════════════════════════════════════════════════
% ЧАСТЬ XVI. РАБОЧАЯ ТЕТРАДЬ: КОНТРОЛЬНЫЕ ВЫЧИСЛЕНИЯ
% ═══════════════════════════════════════════════════════════════════
\part{Рабочая тетрадь: контрольные вычисления}

\section{Как пользоваться тетрадью}\label{sec:workbook}

Десять контрольных вычислений ниже --- самопроверка читателя и
одновременно готовые лабораторные опыты: каждое имеет полное
решение в тексте и пункт меню лаборатории, где оно воспроизводится
автоматически. Вычисления подобраны так, чтобы покрыть все слои
конвейера: целые числа, фазы, интегралы, решётки, башни.

\subsection*{Задача 1. Род и перепись для $N=7$}
Проверить формулу рода для кривой Ферма $N=7$ и построить перепись
характеров по проводникам.

\emph{Решение.} Для $C_7=\{x^7+y^7=z^7\}$ формула рода даёт
$g=\tfrac{6\cdot5}2=15$. Перепись по формуле~\eqref{eq:mobius}:
единственный делитель уровня $7$ --- простое число, поэтому
$h_7=\sum_{m\mid7}\mu(m)\,c(7/m)=c(7)-c(1)
=\tfrac{6\cdot5}2-0=15$; сумма переписи $1\cdot15=15=g$ ---
согласована. Важное замечание: квартика Клейна
(раздел~\ref{sec:klein}) --- гладкая плоская кривая степени $4$ с
родом $\tfrac{3\cdot2}2=3$ --- является \emph{другой} кривой,
несмотря на присутствие семёрки в её группе симметрий
$\Gamma(2,3,7)$; формула рода применяется к степени плоской
модели, и смешение моделей недопустимо. Пункт лаборатории:
«Конструктор $\to$ перепись для произвольного $N$» (ввод $N=7$
даёт $h_7=15$, сумма $15$).

\subsection*{Задача 2. Тройка и поправки для $N=5$}
Вычислить полный набор поправок для поворота порядка $5$ при
$k\in\{1,22\}$.

\emph{Решение.} $\delta=\pi/5=0.6283185$; $\delta^2/2=0.1973921$;
$\gamma(k{=}1)=\delta^4=0.1558545$; $\delta_{\mathrm{eff}}(k{=}1)
=\delta^5=0.0979263$; $\gamma(k{=}22)=0.0070843$;
$\delta_{\mathrm{eff}}(k{=}22)=0.0044512$.
$b_{Ch}(5)=1-\cos(72^\circ)=1-\tfrac{\sqrt5-1}4
=\tfrac{5-\sqrt5}4\approx0.6909830$. Все числа согласованы со
строкой $N=5$ таблицы приложения~\ref{app:fulldata}. Пункт:
«Стенды $\to$ Тор $\to$ таблица $\gamma,\delta_{\mathrm{eff}}$».

\subsection*{Задача 3. Фаза периода вручную}
Для $N=15$, характера $(a,b)=(2,3)$ и данных обхода $(r,s)=(2,1)$
вычислить фазу периода в точной форме.

\emph{Решение.} Фаза $\zeta_{15}^{ra+sb}=\zeta_{15}^{2\cdot2+1
\cdot3}=\zeta_{15}^{7}$. Угол $\frac{2\pi\cdot7}{15}$ --- точное
значение; период $P=\frac1{15}\zeta_{15}^7\Omega_{2,3}$ с
$\Omega_{2,3}=\Ga(2/15)\Ga(3/15)/\Ga(5/15)$. Проверка лаборатории
(V3): аргумент численного интеграла равен $\frac{14\pi}{15}$ с
отклонением $0.0$. Пункт: «Стенд N=15/30 $\to$ V2/V3».

\subsection*{Задача 4. Пересечение прямых K3}
Вычислить $L_1(1,2)\cdot L_2(3,0)$ по правилам
раздела~\ref{sec:k3}.

\emph{Решение.} Семейства разные ($\{1,2\}$), условие
$a_1+b_2\equiv a_2+b_1\pmod4$: $1+0=1$ против $3+2=5\equiv1$ ---
совпадают, пересечение $1$. Для пары $L_1(1,2)\cdot L_1(1,2)$ ---
самопересечение $-2$. Пункт: «Стенды $\to$ K3» (матрица
пересечений печатается в отчёте).

\subsection*{Задача 5. Кратность характера}
По данным теоремы~\ref{th:C1} вычислить сумму размерностей изотипных
частей и сравнить с $b_2$.

\emph{Решение.} $1+7+7+7=22=b_2(K3)$ --- согласовано. Контроль
через проекторы: $\rank\Pi_{\mathrm{triv}}=1$,
$\rank\Pi_{i}=\rank\Pi_{-1}=\rank\Pi_{-i}=7$. Пункт:
«Сертификаты $\to$ C».

\subsection*{Задача 6. Время завершимости}
Вычислить $t^*$ для сетки $24\times36$ и шага $(3,5)$.

\emph{Решение.} $W/\gcd(a,W)=24/3=8$; $H/\gcd(b,H)=36/1=36$
(так как $\gcd(5,36)=1$); $t^*=\lcm(8,36)=72$. Проверка перебором
лаборатории даёт ровно $72$ перехода до возврата. Пункт:
«Сертификаты $\to$ E $\to$ произвольная сетка».

\subsection*{Задача 7. Радиус единственности}
Доказать, что при $Q=480$ и $100$ цифрах точности интервалы
единственности сертификата~F выполняются с запасом.

\emph{Решение.} Порог единственности $2\varepsilon<1/Q^2
=\tfrac1{480^2}=\tfrac1{230400}\approx4.34\cdot10^{-6}$. Сто цифр
дают $\varepsilon\sim10^{-100}$ --- запас $94$ порядка. Для стенда
все знаменатели $\le480=Q_{\text{стенд}}$
(теорема~\ref{th:H56}), следовательно скан замыкается без
дополнительных измерений. Пункт: «Сертификаты $\to$ F, H».

\subsection*{Задача 8. Спектральная дробь CM}
Вычислить $\lambda_1$ для решётки $\tau=\tfrac{1+\sqrt{-15}}2$.

\emph{Решение.} $\im\tau=\tfrac{\sqrt{15}}2$; приведённая решётка:
$\lambda_1=4\pi^2/(\im\tau)^2=4\pi^2/\tfrac{15}4
=\tfrac{16\pi^2}{15}$ --- семейство $16\pi^2/d$ при $d=15$
(теорема~\ref{th:certG}). Норма кратчайшего вектора
$|1\cdot\tau-0|^2=|(\tfrac12)^2+\tfrac{15}4|=\tfrac{16}4=4$...
точнее: $|\tau|^2=\tfrac14+\tfrac{15}4=4$, и деление на
$(\im\tau)^2$ даёт $\tfrac{4}{15/4}=\tfrac{16}{15}$ ---
$\lambda_1=(2\pi)^2\cdot\tfrac{16}{15}=\tfrac{16\pi^2}{15}$.
Согласовано. Пункт: «Сертификаты $\to$ G $\to$ решётки».

\subsection*{Задача 9. SNF-произведение}
Вычислить произведение инвариантных множителей стенда и сверить с
дискриминантом.

\emph{Решение.} $(1,1,5,5,15,15,15,15)$: произведение
$5\cdot5\cdot15^4=25\cdot50625=1265625=3^4\cdot5^6$;
$\sqrt{1265625}=1125$ --- целое. Совпадение с
теоремой~\ref{th:H1} точное. Пункт: «Сертификаты $\to$ H1».

\subsection*{Задача 10. Радикал $b_{Ch}(15)$}
Вывести радикал~\eqref{eq:b15} из формулы косинуса разности.

\emph{Решение.} $\cos24^\circ=\cos(60^\circ-36^\circ)
=\cos60\cos36+\sin60\sin36=\tfrac12\cdot\tfrac{1+\sqrt5}4
+\tfrac{\sqrt3}2\cdot\tfrac{\sqrt{10-2\sqrt5}}4
=\tfrac{1+\sqrt5+\sqrt3\sqrt{10-2\sqrt5}}8$.
Тогда $b_{Ch}(15)=1-\cos24^\circ$ --- формула~\eqref{eq:b15}.
Число: $0.0864539\ldots$; проверка Lean: выражение
$\tfrac{1+\sqrt5+\sqrt3\sqrt{10-2\sqrt5}}8$ сверено с
$1-\cos(2\pi/15)$ численно в ядре. Пункты: «Теорема 2»,
«Lean-панель».

\section{Ответы на типовые вопросы воспроизведения}\label{sec:faq}

\emph{Вопрос: почему V7 не использует PSLQ?}
Ответ: полнота --- свойство ранга, а не соотношений. Ранг матрицы
функционалов $=g$ означает ровно столько независимых периодов,
сколько форм; поиск соотношений не нужен и не производится. PSLQ
используется только в служебной роли --- как независимое
восстановление уже выведенных показателей (H2), с точной
аналитической сертификацией.

\emph{Вопрос: почему соотношения Римана не проверяются численно на
$91\times182$?}
Ответ: они доказаны тождественно (теорема~\ref{th:riemann});
численная проверка тождества с точностью $10^{-35}$ была бы
\emph{слабее} доказательства. Числа остаются одноразовой
санити-строкой вне текста.

\emph{Вопрос: чем сертификат отличается от теста?}
Ответ: тест сверяет код с самим собой; сертификат сверяет
реализацию с независимым методом (интегрирование без
$\Ga$-функций, целочисленная арифметика, ядро Lean). Правило
независимости (раздел~\ref{sec:rules}) запрещает сертификату
использовать тот же алгоритм, что и проверяемое вычисление.

\emph{Вопрос: что нужно, чтобы добавить новую ступень лестницы?}
Ответ: только тройку $(n,B,\lambda_0)$ и перепись характеров
новой геометрии. Конвейер, протокол, Lean-панель и мультиязычная
верификация работают без изменений; конструктор экспериментов
(меню лаборатории) принимает произвольные параметры. Стенд
$N=15/30$ демонстрирует полный цикл такого подключения.
